\indent 
    
    The GMES School is a comprehensive system designed to serve educational institutions from kindergarten to high school. The system helps schools manage their daily operations with easy-to-use tools. The system is built based on the four core modules:
    \begin{itemize}[itemsep=-0.1cm, topsep=0.1cm]
        \item \textbf{Personnel management}: This module focuses on managing information such as personal details, contracts, education qualifications, and relationships of teachers, staff, and students. This module is fundamental to the system because all other features rely on this module's data.
        \item \textbf{Attedance management}: This module is used to track the time in and out of personnel through an attendance machine. By using these records we can calculate the work hours, and salary, and generate attendance reports of people.
        \item \textbf{Leave request management}: This module focuses on creating, managing, and processing leave requests (approve, decline) of people. This module will help the school manage leave requests easily.
        \item \textbf{Billing management}: This module provides financial management solutions for educational institutions such as handles the payment systems for staff salaries and student fees.
    \end{itemize}
    Besides the four core modules, the system provides several supporting modules that are equally important for creating a comprehensive, differentiated school management system. These supporting modules integrates with the core modules to form a unified ecosystem for educational institutions. These supporting modules include:
    \begin{itemize} [itemsep=-0.1cm, topsep=0.1cm]
        \item \textbf{Authentication management}: This module controls system access through secure login/logout processes and role-based authorization, ensuring users only access permitted features and data.
        \item \textbf{Schedule management}: This module comprehensively manages work schedules and academic timetables, enabling flexible task allocation and eliminating scheduling conflicts.
        \item \textbf{Device management}: This module is responsible for the management and control all attendance machines, allowing remote monitoring, configuration, and data synchronization across multiple school institutions.
        \item \textbf{Organization structure management}: This module maintains the school's hierarchy through a tree-like system, enabling administrators to manage relationships from the institution level down to individual departments, grades, and classes.
        \item \textbf{Bus management}: The module handles all aspects related to school transportation services within educational institutions. It also integrates with attendance modules to ensure student safety while optimizing transport operations.
        \item \textbf{Newsfeed management}: This module facilitates school communication by allowing teachers to share announcements, events, and educational content with parents through multimedia posts, keeping families actively engaged and informed about their children's school experience.
        \item \textbf{Assessment management}: This module handles performance reviews for staff and academic evaluations for students, ensuring consistent standards and providing valuable insights for improvement.
        \item \textbf{Health management}: The module tracks student health information and medication schedules, ensuring student well-being, supporting medical compliance, and giving parents peace of mind about their child's healthcare at school.
        \item \textbf{Notification management}: This module is responsible for keeping all stakeholders, including employees, teachers, and parents informed about critical events, system updates, and important announcements.
    \end{itemize}
    In the GMES School management system, five distinct roles are defined to align with their specific duties and organizational hierarchies. These roles ensure users can only access features and information appropriate to their position and responsibilities within the educational institution. These roles are: 
    \begin{itemize}[itemsep=-0.1cm, topsep=0.1cm]
        \item \textbf{School administrator}: A person in the education institutions who controls entire system operations for their specific school.  
        \item \textbf{Staff}: Office workers or kitchen staff who handle paperwork or prepare meals.
        \item \textbf{Teacher}: A person with enough academic knowledge to teach students.
        \item \textbf{Transportation staff}: A person responsible for driving the bus or managing the bus.
        \item \textbf{Parents}: Parents of students who study at school.
    \end{itemize}